\chapter{Changes in sweepgen and waLBerla}\label{r03:ch:walberla}

\emph{AA streaming and kernels without output in sweepgen; the halo exchange for AA streaming
and its faster rewrite in waLBerla's v8 core; the benchmark application; the direction-split
kernel generated from lbmpy; and the build flags.}

The fork is \href{\rWLurl/tree/\rWLsha}{\texttt{mirrored-walberla}}, branch
\texttt{rvv-support}, commit \rcommitwl. The changes described here touch 14 files
(\href{\rWLurl/compare/b78f2d8791589fa38d0e407a3447865cb8ce227f...\rWLsha}{\texttt{b78f2d87..efa2b70c}}:
1020 lines added, 93 removed).

\section{sweepgen: AA streaming and kernels without output}\label{r03:sec:wl-sweepgen}

Both changes are in the class \code{LbmBulk} in
\ghwl{sweepgen/src/sweepgen/prefabs/lbm_bulk.py}, which generates the LBM sweeps of a method
(\cref{r03:sec:bg-stack}).
\begin{description}
  \item[AA streaming.] When lbmpy's configuration asks for it
        (\code{LBMConfig(streaming_pattern="aa")}), \code{LbmBulk} builds lbmpy's update rules
        for the even and the odd step (\code{Timestep.EVEN} and \code{Timestep.ODD}). It
        generates two sweeps, \code{StreamCollideEven} and \code{StreamCollideOdd}, that work on
        one PDF field: no second field and no swap. \code{InitPdfs} writes the PDFs where the
        even step will read them, so each cell writes only its own slots. Pull streaming is
        unchanged.
  \item[Kernels without output.] With \code{output_free_sweeps=True}, \code{LbmBulk} also
        generates sweeps that update the PDFs but do not store density and velocity:
        \begin{itemize}
          \item pull: \code{StreamCollideNoOutput};
          \item AA: \code{StreamCollideEvenNoOutput} and \code{StreamCollideOddNoOutput}.
        \end{itemize}
        They come from lbmpy update rules without output fields. The PDF update is the same,
        and there are 4 fewer store streams. For pull, sweepgen swaps the field with a
        temporary buffer that belongs to the sweep object; the extra sweep object therefore has
        a buffer of its own. Alternating the two sweeps stays correct, at the cost of one more
        PDF-sized buffer.
\end{description}

\section{The halo exchange in the v8 core}\label{r03:sec:wl-halo}

\paragraph{Pack infos for AA streaming.} AA streaming needs a different halo exchange from
pull streaming, because the odd step reads from and writes into the ghost layer. Two new pack
infos (\cref{r03:sec:bg-blocks}) and functions that create them were added to
\ghwl{include/walberla/v8/halo_exchange/fields/GenericFieldPackInfos.hpp} and
\ghwl{include/walberla/v8/halo_exchange/fields/PackInfoSelection.hpp}:
\begin{description}
  \item[\code{GenericAAPullPackInfo}, created by \code{halo_exchange::aaPullSync<Stencil>}] runs
        after even steps. The odd step reads slot $\bar\imath$ at $x - c_i$, so the ghost layer
        must hold the slots of the inverse directions of those the pull exchange
        (\code{streamPullSync}) sends.
  \item[\code{GenericAAPushPackInfo}, created by \code{halo_exchange::aaPushSync<Stencil>}] runs
        after odd steps. The odd step of a border cell writes slot $i$ at $x + c_i$, which can be
        a ghost cell. This exchange sends those values back to the block that owns that cell.
        At the edges of a face, some ghost slots were not written by the sending block, because
        the cell that would write them, $g - c_i$ for ghost cell $g$, lies outside it. The owner
        receives those values from another neighbour, across an edge. The exchange must
        therefore skip them; otherwise the order of the unpack operations would decide whether a
        stale value overwrites the correct one.
\end{description}
Slot indices are computed with waLBerla's stencil functions \code{Stencil::idx} and
\code{Stencil::invDirIdx}.

\paragraph{One slot at a time.} A profile of the f32 AA application showed the two AA pack
infos' \code{localCopy} at 17.8\,\% of all samples, against 68.8\,\% for the kernels: 21\,\% of
a time step (\cref{r03:tab:perf}). With one block, every direction of the halo exchange is a
periodic copy of the block onto itself (\cref{r03:sec:bg-blocks}). The copy looped over the
cells of a face and, for each cell, over the 5 slots that cross that face. In the fzyx layout
these 5 slots are in 5 different slices, so on an $x$ face every copied value sat on its own
cache line. The AA push exchange also tested for every value whether its slot had been written.
The pull pack info and the two AA pack infos now move one PDF slot at a time
(\cref{r03:fig:halo}):
\begin{itemize}
  \item \code{localCopy}, \code{pack} and \code{unpack} loop over the slots in the outer loop and
        over the cells in the inner loop (helper functions \code{detail::copySlot},
        \code{packSlot}, \code{unpackSlot}). In the fzyx layout each pass reads and writes one
        slice in storage order.
  \item The MPI message holds one contiguous segment per slot instead of one per cell. Sender
        and receiver use the same class, so this layout is internal; the message size
        (\code{packetSize}) is unchanged.
  \item The AA push exchange computes the cells to copy once per slot, as a box: the interior
        shifted by $c_i$ and cut to the ghost layer (function \code{writtenGhosts}).
  \item Execution on GPUs takes the same path, with one kernel launch per slot instead of one per
        direction.
\end{itemize}
The rewrite passes the correctness checks (\cref{r03:sec:ver-wl}); its effect on the
application's time step was still being measured (\cref{r03:ch:open}).

\begin{figure}[htbp]
  \centering
  \resizebox{\linewidth}{!}{% localCopy on an x face, direction W: the 5 PDF slots with c_x = -1 (D3Q19 slots 3, 7, 9,
% 13, 17), each in its own slice of the
% fzyx field. Before: for every cell, all 5 slots (the access order jumps between slices). After: for
% every slot, all cells of the face (one slice at a time). Numbers = order of the copies.
\begin{tikzpicture}[every node/.style={font=\scriptsize}]
  \def\rthreeSlices#1#2#3{% x offset, title, order macro (1 = cell-major, 2 = slot-major)
    \node[font=\footnotesize\bfseries,anchor=west] at (#1,1.55) {#2};
    \foreach \q/\slotname in {0/3 (W),1/7 (NW),2/9 (SW),3/13 (TW),4/17 (BW)} {
      \pgfmathsetmacro{\rthreeSX}{#1+\q*1.45}
      \node[r03label] at (\rthreeSX+0.5,1.15) {\slotname};
      \foreach \c in {0,...,3} {
        \pgfmathsetmacro{\rthreeSY}{0.75-\c*0.42}
        \draw[psgray!50,line width=0.2pt] (\rthreeSX,\rthreeSY) rectangle ++(1.0,-0.42);
        \fill[psblue!14] (\rthreeSX+0.62,\rthreeSY) rectangle ++(0.38,-0.42);
        \ifnum#3=1 \pgfmathtruncatemacro{\rthreeOrd}{\c*5+\q+1}\else\pgfmathtruncatemacro{\rthreeOrd}{\q*4+\c+1}\fi
        \node[font=\tiny\bfseries,text=psblue!80!black] at (\rthreeSX+0.81,\rthreeSY-0.21) {\rthreeOrd};
      }
    }
  }
  \rthreeSlices{0}{before: cells outside, slots inside}{1}
  \rthreeSlices{7.6}{after: slots outside, cells inside}{2}
  \node[r03label,anchor=north west,text width=146mm] at (0,-1.1)
    {Columns: the 5 slots with $c_x = -1$ (slot index and direction). Each small grid is one slot's slice of the field, seen on the face (rows: 4 of the 64 cells along
     $y$); shaded: the column copied. Before, consecutive copies hop between 5 slices, and the AA push
     variant tests per value whether the slot was written. After, each pass walks one slice; the AA
     push variant copies a precomputed box per slot (\texttt{writtenGhosts}). On $y$ and $z$ faces the
     copied runs are contiguous rows of 64 cells.};
\end{tikzpicture}
}
  \caption[Halo copy: cell by cell and slot by slot]{\texttt{localCopy} on an $x$ face, before
    and after the rewrite. Numbers give the order of the copies.}\label{r03:fig:halo}
\end{figure}

\section{The benchmark application}\label{r03:sec:wl-app}

\ghwl{apps/benchmarks/SweepgenLbmCPU/SweepgenLbmCPU.cpp},
\ghwl{apps/benchmarks/SweepgenLbmCPU/SweepgenLbmCPU.py} (the Python script that generates the
application's kernels) and \ghwl{apps/benchmarks/SweepgenLbmCPU/CMakeLists.txt}:
\begin{itemize}
  \item New options of the generator script:
        \begin{itemize}
          \item \code{--prefetch-distance N}, for the RVV variants;
          \item \texttt{-{}-streaming-pattern \{pull,aa\}};
          \item \code{--split BY,GM,GP}, the direction-split kernel (\cref{r03:sec:wl-split}).
        \end{itemize}
        The script always asks \code{LbmBulk} for the kernels without output as well.
  \item The CMake list of variants is built by a helper function; the variants are listed in
        \cref{r03:tab:variants}. The compile definitions \code{SWEEPGEN_LBM_AA} and
        \code{SWEEPGEN_LBM_SPLIT} select the time loop.
  \item AA time loop: even and odd steps alternate, \code{aaPullSync} after even and
        \code{aaPushSync} after odd steps (\cref{r03:fig:timestep}).
  \item The environment variable \code{SWEEPGEN_LBM_OUTPUT=last} makes every step except the
        last use the kernel without output. The default, \code{every}, writes the output in
        every step, as before. The result line prints \code{output=every} or
        \code{output=last}.
  \item \textbf{Correctness fix.} The pull benchmark as it was before this work ran one halo
        exchange before the first time step. \code{InitPdfs} places the PDF of direction $i$ of
        cell $x$ in slot $i$ of cell $x - c_i$; near a block border that is a ghost cell, and
        these are exactly the values the first pull step reads. The exchange overwrote them with
        border slots that nothing had written yet, so the result depended on how the domain was
        divided into blocks: after one step, 382.26 with one block of $64^3$ and 365.71 with 8
        blocks of $32^3$, where 399.23 is correct. The exchange was removed. Checksums of runs
        before the fix agree between variants, but are not the physically correct values; the
        timed loop is unchanged, so speeds are unaffected.
\end{itemize}

\section{The direction-split kernel, generated from lbmpy}\label{r03:sec:wl-split}

\ghwl{apps/benchmarks/SweepgenLbmCPU/lbm_split.py} (new) generates the structure of the hand
kernel of \cref{r03:sec:lever-aa} from lbmpy's description of the method. The lab's generator
\ghws{misc/x60-lab/lab/gensplit.py} imports the same module, so lab and application use one
implementation. Three parameters choose the structure: BY rows per block, GM directions per
moment pass, and GP pairs of opposite directions per collision pass.

\begin{figure}[htbp]
  \centering
  % Passes of the direction-split AA stream/collide for one block of BY rows of one z-plane
% (lbm_split.py, BY=2 GM=5 GP=3 as in the f32 build). Used in chapter 06.
\begin{tikzpicture}[node distance=3mm and 4mm]
  \node[r03pass,text width=17mm] (m0) {m0\\[-1pt]\scriptsize dirs 0--4};
  \node[r03pass,text width=17mm,right=of m0] (m1) {m1\\[-1pt]\scriptsize dirs 5--9};
  \node[r03pass,text width=17mm,right=of m1] (m2) {m2\\[-1pt]\scriptsize dirs 10--14};
  \node[r03pass,text width=17mm,right=of m2] (m3) {m3\\[-1pt]\scriptsize dirs 15--18};
  \node[r03box,text width=20mm,below=9mm of m1.south east,anchor=north] (acc)
    {\texttt{acc}\\[-1pt]\scriptsize $\delta\rho$, $j_x$, $j_y$, $j_z$};
  \node[r03pass,text width=20mm,right=12mm of acc] (macro) {macro\\[-1pt]\scriptsize $\rho$, $u$ (+ outputs)};
  \node[r03box,text width=24mm,right=8mm of macro] (scr)
    {\texttt{scr}\\[-1pt]\scriptsize $\delta\rho$, $\rho$, $u_x$, $u_y$, $u_z$};
  \node[r03pass,text width=20mm,below=9mm of acc] (c0) {c0\\[-1pt]\scriptsize rest + 3 pairs};
  \node[r03pass,text width=20mm,right=of c0] (c1) {c1\\[-1pt]\scriptsize 3 pairs};
  \node[r03pass,text width=20mm,right=of c1] (c2) {c2\\[-1pt]\scriptsize 3 pairs};
  \foreach \m in {m0,m1,m2,m3} \draw[r03arrow] (\m.south) -- (acc.north);
  \draw[r03arrow] (acc) -- (macro);
  \draw[r03arrow] (macro) -- (scr);
  \foreach \c in {c0,c1,c2} \draw[r03arrow] (scr.south) -- (\c.north);
  \node[r03label,above=1mm of m1.north east] {moment passes: read the PDFs of a few directions, accumulate};
  \node[r03label,below=1mm of c1.south] {collision passes: reload the PDFs (L1/L2), relax, store in place};
\end{tikzpicture}

  \caption[Passes of the direction-split AA kernel]{Passes of the direction-split AA kernel for
    one block of BY rows of one $z$-plane, with GM~=~5 directions per moment pass and GP~=~3
    pairs per collision pass. Each pass is one kernel generated by \ps{}. \texttt{acc} and
    \texttt{scr} are scratch arrays of the size of one block.}\label{r03:fig:split}
\end{figure}

\paragraph{What is generated.} The class \code{SplitAA} builds, for the even and for the odd
step:
\begin{itemize}
  \item \textbf{moment passes}: for a group of directions, load the PDFs with lbmpy's AA read
        positions and add $\delta\rho = \sum f_i$ and the momentum $j = \sum c_i f_i$ into the
        scratch array \code{acc} (the first pass assigns, the others add);
  \item \textbf{a macroscopic pass}: $\rho = \delta\rho + 1$ and $u = j / \rho$, into the scratch
        array \code{scr}, and in one of two versions also into the output fields;
  \item \textbf{collision passes}: for a group of opposite pairs (the rest direction joins the
        first group), load all PDFs of the group, then store $f_i + \omega\,(f^{\mathrm{eq}}_i -
        f_i)$ at lbmpy's AA write positions.
\end{itemize}
Colliding each direction separately is only valid for SRT, where every PDF relaxes on its own;
the generator rejects methods with more than one relaxation rate.

\paragraph{The equilibrium comes from lbmpy, and is checked.} A first version copied the hand
kernel's textbook polynomial $t = 3\,c\cdot u + \tfrac92 (c\cdot u)^2 - \tfrac32 u^2$. Its
symbolic comparison with lbmpy's equilibrium (\code{method.get_equilibrium_terms()}) failed.
lbmpy's D3Q19 SRT equilibrium differs from the textbook one: for the rest direction it is
$w_0(\delta\rho - \rho u^2)$ instead of $w_0(\delta\rho - \tfrac32 \rho u^2)$, and for face and
edge directions it differs by multiples of $u_a^2$, the squared velocity components. The
difference carries no mass, momentum or second moment, so it only affects the
\term{non-hydrodynamic moments}: combinations of PDFs that do not enter the fluid equations. The
generator now:
\begin{enumerate}
  \item takes $t_i = (f^{\mathrm{eq}}_i / w_i - \delta\rho) / \rho$ from lbmpy and checks
        symbolically that $f^{\mathrm{eq}}_i = w_i (\delta\rho + \rho\, t_i)$ for every $i$;
  \item splits each opposite pair $(i, j)$ into an even part $e = (t_i + t_j)/2$ and an odd part
        $o = (t_i - t_j)/2$, so that $t_i = e + o$ and $t_j = e - o$: one $e$ and one $o$ per
        pair;
  \item lets sympy find common sub-expressions within each pass.
\end{enumerate}
The hand kernels therefore compute a slightly different model. Under QEMU (f64, $16^3$ cells,
20 steps) their $|u|^2$ checksum differs from that of the generated kernels by
$6\cdot10^{-11}$ relative, and their weighted PDF sum, which sees the non-hydrodynamic moments,
by 1.3\,\%. This does not affect their speed.

\paragraph{Driver.} \code{SplitAA.driver_code} emits a C++ function with the parameters
\code{(f, rho, u, n0, n1, n2, strides..., odd, output, omega)}. For each $z$-plane and each block
of up to BY rows, it moves the field pointers to the block and calls the passes with the
iteration space $(n_0, \text{rows}, 1)$. The scratch arrays are one \code{static std::vector}
sized for one block, each array aligned to 64~bytes. The versions with and without output share
all passes except the macroscopic one, and each version keeps its own order of passes.

\paragraph{In the application.} \code{SweepgenLbmCPU.py --split BY,GM,GP} (RVV variant with AA
streaming only) generates the passes for the application's vector target, without ghost layers
of their own (\code{ghost_layers = 0}), and the driver function \code{LBM_SPLIT::streamCollide}.
The variant \code{explicit_cse_aa_split} uses \code{2,5,3} in single and \code{1,5,3} in double
precision, the best settings of the hand kernel (\cref{r03:sec:lever-aa}), and defines
\code{SWEEPGEN_LBM_SPLIT}. The application then calls the driver instead of the AA sweeps; halo
exchange and initialisation are those of the AA variant.

\section{Build flags: the X60 scheduling model}\label{r03:sec:wl-cmake}

\ghwl{cmake/compileroptions/GNU.cmake} and \ghwl{cmake/compileroptions/Clang.cmake}, in the RISC-V
branch of the option \code{WALBERLA_OPTIMIZE_FOR_LOCALHOST} (optimise for the build machine):
\begin{itemize}
  \item read the core type from \file{/proc/cpuinfo} (\code{uarch : spacemit,x60}
        $\rightarrow$ \code{-mtune=spacemit-x60});
  \item add the flag only if the compiler accepts it (CMake's \code{check_cxx_compiler_flag}).
\end{itemize}
Neither compiler derives the model from \code{-mcpu=native}: Clang~18 resolves it to a generic
core. GCC~14 rejects the flag, so GCC builds are unchanged.

\section{Tests}\label{r03:sec:wl-tests}

\ghwl{tests/sweepgen/LbmScenarios/AAMatchesPull.hpp} (new). waLBerla's LBM scenario tests use a
TRT method with a body force (lbmpy's \code{ForceModel.HE}), generated in
\ghwl{tests/sweepgen/LbmScenarios/LbmAlgorithms.py}. That script now also generates the method
with AA streaming, as \code{LBM_AA}. The new test runs pull and AA side by side on a periodic flow
with a sinusoidal velocity along every axis. After each of 9 time steps it compares $\rho$ and
$u$ cell by cell (absolute tolerance $10^{-14}$, relative $10^{-12}$). CTest, CMake's test runner,
runs it with 1, 2 and 4 MPI processes. It therefore covers both the periodic copy of one block
onto itself (\code{localCopy}) and the exchange between blocks (\code{pack}/\code{unpack}).

\section{Scripts in the workspace}\label{r03:sec:wl-scripts}

\ghws{scripts/walberla_lbm_bench.sh} builds and runs the AA and prefetch variants by default;
\code{TARGETS} overrides the list. \ghws{scripts/walberla_bench_mpi.sh} (several MPI processes)
takes \code{CONFIGS}, \code{VARIANTS} and \code{LOGDIR}. \ghws{scripts/walberla_bench_summary.py}
prints summary tables and knows the new variants. The working tree also adds the split variant to
these default lists; that edit is newer than \rcommitws.
